\chapter{Sobolev 空间}

\begin{solution}{sobolev-absolute-value}
\(|x|\in W^{0,p}(-1,1)=L^p(-1,1)\) 对所有 \(1\le p\le\infty\) 成立，一阶弱导数
\(\operatorname{sgn}x\in L^\infty\)，故属于所有 \(W^{1,p}\)。
二阶分布导数为 \(2\delta_0\)，不由局部可积函数表示，所以不属于任何
\(W^{2,p}\)，更不属于更高阶空间。
\end{solution}

\begin{solution}{sobolev-heaviside-failure}
若 \(H\in W^{1,p}\)，其弱导数应由某个 \(g\in L^p\subset L^1_{\rm loc}\)
表示。但分布计算给 \(DH=\delta_0\)，而 Dirac 质量不是正则分布，矛盾。
\end{solution}

\begin{solution}{sobolev-weak-derivative-limit}
对 \(\varphi\in C_c^\infty\)，由 Hölder 不等式，
\[
 \int u_n\partial_j\varphi\to\int u\partial_j\varphi,\qquad
 \int(\partial_ju_n)\varphi\to\int g\varphi.
\]
在
\(\int u_n\partial_j\varphi=-\int(\partial_ju_n)\varphi\) 中取极限，得到
\(g\) 是 \(u\) 的弱导数。
\end{solution}

\begin{solution}{sobolev-hilbert-inner-product}
每项 \(L^2\) 内积均满足共轭对称和线性，有限和仍满足。若
\(\langle u,u\rangle_{H^k}=0\)，其中零阶项给 \(\|u\|_2=0\)，故 \(u=0\)
作为等价类。范数平方恰为
\(\sum_{|\alpha|\le k}\|\partial^\alpha u\|_2^2\)，与定义一致。
\end{solution}

\begin{solution}{sobolev-reflexive}
映射
\[
 Ju=(\partial^\alpha u)_{|\alpha|\le k}
\]
把 \(W^{k,p}\) 等距嵌入有限积 \(Y=\prod L^p\)。完备性证明已经表明其像闭。
\(1<p<\infty\) 时每个 \(L^p\) 反射，有限积反射，反射空间的闭子空间反射。
因此 \(W^{k,p}\) 反射。
\end{solution}

\begin{solution}{sobolev-fourier-hone}
若 \(u\in H^1\)，先用梯度为 \(O(R^{-1})\) 的光滑截断 \(\chi_R\) 乘以
\(u\)，再磨光。乘积法则和 \(L^2\) 尾积分表明所得
\(u_n\in C_c^\infty\) 在 \(H^1\) 中逼近 \(u\)。对每个 \(n\)，
\(\widehat{\partial_ju_n}=2\pi i\xi_j\widehat u_n\)。由 Plancherel，
\(\widehat u_n\to\widehat u\)，且 \(\widehat{\partial_ju_n}\to
\widehat{\partial_ju}\) 于 \(L^2\)。对 \(u_n\) 及全部一阶导数的变换同时
抽取几乎处处收敛子列，乘子等式便在该满测集上保留。
因而 \(\widehat{\partial_ju}=2\pi i\xi_j\widehat u\) 几乎处处，
范数恒等式再由等距性得到。不能直接把乘以 \(\xi_j\) 当成连续的 \(L^2\) 算子。

反向令 \(g_j=\mathcal F^{-1}(2\pi i\xi_j\widehat u)\)。对试验函数使用
Parseval，并注意
\(\widehat{\partial_j\varphi}=2\pi i\xi_j\widehat\varphi\)，可得
\(\int u\,\partial_j\varphi=-\int g_j\varphi\)；复值情形按双线性分布配对或
先分实虚部处理。故 \(g_j\) 是弱导数。
\end{solution}

\begin{solution}{sobolev-poincare-interval}
若 \(u\in C_c^\infty(0,L)\)，则
\[
 |u(x)|=\left|\int_0^xu'(t)dt\right|
 \le\sqrt{x}\|u'\|_2.
\]
积分给 \(\|u\|_2\le L/\sqrt2\,\|u'\|_2\)，并延拓到 \(H_0^1\)。
为证明下界，令 \(s(x)=\sin(\pi x/L)\)。先核实 \(s\in H_0^1(0,L)\)。
取光滑截断 \(\chi_\varepsilon\)，在 \([2\varepsilon,L-2\varepsilon]\)
上为一，在端点各长 \(\varepsilon\) 的邻域内为零，满足
\(0\le\chi_\varepsilon\le1\)、\(|\chi_\varepsilon'|\le C/\varepsilon\)。
则 \(\chi_\varepsilon s\in C_c^\infty(0,L)\)。误差仅在总长度
\(4\varepsilon\) 的端点区域上出现，其中 \(|s(x)|\le C\varepsilon\)、
\(|s'(x)|\le C\)，故
\[
 \|(1-\chi_\varepsilon)s\|_2=O(\varepsilon^{3/2}),\qquad
 \|(1-\chi_\varepsilon)s'-\chi_\varepsilon's\|_2=O(\varepsilon^{1/2}).
\]
所以确有 \(H^1\) 逼近。直接积分给
\(\|s\|_2^2=L/2\)、\(\|s'\|_2^2=\pi^2/(2L)\)，任何有效常数因此至少为 \(L/\pi\)。

还能在不预用正弦展开的情况下证明这就是最佳常数。令 \(\lambda=\pi/L\)。
对 \(u\in C_c^\infty(0,L)\)，因 \(s>0\) 于内部，\(v=u/s\) 光滑且紧支撑。
展开 \(u'=s'v+sv'\)，再将交叉项分部积分，得到
\[
 \begin{aligned}
 \int_0^L|u'|^2
 &=\int_0^L\bigl(s^2|v'|^2+(s')^2|v|^2+ss'(|v|^2)'\bigr)\\
 &=\int_0^Ls^2|v'|^2-\int_0^Lss''|v|^2
 =\int_0^Ls^2|v'|^2+\lambda^2\int_0^L|u|^2.
 \end{aligned}
\]
边界项因 \(v\) 紧支撑消失，最后用 \(s''=-\lambda^2s\)。右侧第一项非负，
故 \(\|u\|_2\le (L/\pi)\|u'\|_2\)。由密度延拓到 \(H_0^1\)，结合上述下界即得最佳性。
\end{solution}

\begin{solution}{sobolev-hzero-gradient-norm}
若 \(\|\nabla u\|_2=0\)，Poincaré 不等式给 \(\|u\|_2=0\)，故这是范数。
显然 \(\|\nabla u\|_2\le\|u\|_{H^1}\)，而
\[
 \|u\|_{H^1}^2=\|u\|_2^2+\|\nabla u\|_2^2
 \le(C_P^2+1)\|\nabla u\|_2^2.
\]
所以两范数等价。
\end{solution}

\begin{solution}{sobolev-lax-milgram-complex}
设 \(a\) 第一变量线性、第二变量共轭线性，且
\(|a(u,v)|\le M\|u\|\|v\|\)、
\(\operatorname{Re}a(u,u)\ge\alpha\|u\|^2\)。
用复 Riesz 表示定义 \(A\) 使 \(a(u,v)=\langle Au,v\rangle\)。
则
\[
 \alpha\|u\|^2\le\operatorname{Re}\langle Au,u\rangle
 \le\|Au\|\|u\|,
\]
故 \(A\) 有下界、单射且像闭。若 \(w\perp\operatorname{ran}A\)，取
\(u=w\) 得 \(a(w,w)=0\)，与强制性矛盾，故像稠密并等于全空间。对连续
共轭线性右端 \(F(v)=\langle f,v\rangle\) 解 \(Au=f\)，估计同实情形。
\end{solution}

\begin{solution}{sobolev-poisson-energy}
若 \(u\) 是弱解，对任意 \(w\in H_0^1\)，
\[
 J(u+w)-J(u)
 =\int\nabla u\cdot\nabla w-\int fw+\frac12\int|\nabla w|^2
 =\frac12\|\nabla w\|_2^2\ge0.
\]
等号由 Poincaré 不等式只在 \(w=0\) 时成立，故极小点唯一。反之若 \(u\)
极小，对 \(J(u+tw)\) 在 \(t=0\) 求导，得到弱方程。
\end{solution}

\begin{solution}{sobolev-elliptic-coefficients}
令
\[
 a(u,v)=\int_\Omega A(x)\nabla u\cdot\nabla v,\qquad
 F(v)=\int_\Omega fv.
\]
矩阵有界给 \(|a(u,v)|\le C\|\nabla u\|_2\|\nabla v\|_2\)，椭圆性给
\(a(u,u)\ge\lambda\|\nabla u\|_2^2\)。Poincaré 不等式给
\(|F(v)|\le C_P\|f\|_2\|\nabla v\|_2\)。在以梯度为范数的 \(H_0^1\) 上
应用 Lax--Milgram，得到唯一弱解及
\(\|\nabla u\|_2\le C_P\lambda^{-1}\|f\|_2\)。
\end{solution}

\begin{solution}{sobolev-lower-order-term}
因为 \(\Omega\) 有界，Poincaré 不等式可用。先核对有界性：
\[
 \left|\int c uv\right|
 \le\|c\|_\infty\|u\|_2\|v\|_2
 \le C_P^2\|c\|_\infty\|\nabla u\|_2\|\nabla v\|_2.
\]
加上梯度项后，双线性型在梯度范数下有界。
\(c\ge0\) 时
\[
 a(u,u)=\|\nabla u\|_2^2+\int c|u|^2\ge\|\nabla u\|_2^2,
\]
故强制。若 \(c_-= \max\{-c,0\}\in L^\infty\)，则
\[
 a(u,u)\ge\|\nabla u\|_2^2-\|c_-\|_\infty\|u\|_2^2
 \ge(1-C_P^2\|c_-\|_\infty)\|\nabla u\|_2^2.
\]
因此 \(C_P^2\|c_-\|_\infty<1\) 是一个充分的小负部条件。
\end{solution}

\begin{solution}{sobolev-weak-minimization}
取极小化序列 \(u_n\)。强制性与有界线性项给
\[
 J(u_n)\ge\frac\alpha2\|u_n\|^2-\|F\|\|u_n\|,
\]
故 \(u_n\) 有界。Hilbert 空间反射，抽取 \(u_{n_j}\rightharpoonup u\)。
连续线性泛函 \(F\) 弱连续；由 \(a\) 对称正定所产生的等价范数弱下半连续，
故 \(J(u)\le\liminf J(u_{n_j})\)，得到极小点。严格凸性给唯一。对一阶变分
求导得到 \(a(u,v)=F(v)\)，与 Lax--Milgram 解一致。此法突出紧性和下半
连续性，算子证明则不要求对称并直接给出可逆性。
\end{solution}

\begin{solution}{sobolev-one-d-continuous-representative}
\chapterref{chap:distributions-weak-derivatives}的一维刻画给一个局部绝对连续代表 \(U\)。固定 \(a_0\in(a,b)\)，有
\[
 U(x)=U(a_0)+\int_{a_0}^{x}u'(t)\,dt.
\]
由于区间长度有限，\(u'\in L^p(a,b)\subset L^1(a,b)\)，右侧在两端均有有限
极限，按同一积分公式定义 \(U(a),U(b)\)。积分函数的绝对连续性证明
\(U\in AC[a,b]\)。若另有连续代表，两者在内部几乎处处相同，连续性使它们
在内部处处相同，端点极限也相同，故此代表唯一。

对 \(a\le x<y\le b\)，若 \(1<p<\infty\)，Hölder 不等式给
\[
 |U(x)-U(y)|
 \le\int_x^y|u'(t)|dt
 \le\|u'\|_{L^p(a,b)}|x-y|^{1-1/p}.
\]
交换 \(x,y\) 得另一方向。\(p=\infty\) 时直接用
\(\int_x^y|u'|\le\|u'\|_\infty(y-x)\)，得到 Lipschitz 估计。

\(p=1\) 的反例可在 \([0,1]\) 上取
\[
 U(0)=0,\qquad U(x)=\frac1{\log(e/x)}\quad(0<x\le1).
\]
其导数 \(g(x)=1/[x\log^2(e/x)]\) 满足 \(\int_0^1g=1\)，且
\(U(x)=\int_0^xg\)，故 \(U\in AC[0,1]\subset W^{1,1}(0,1)\)。
对每个 \(\alpha>0\)，\(U(x)/x^\alpha\to\infty\)（令 \(t=\log(e/x)\)，
此比值为 \(e^{\alpha(t-1)}/t\)），所以它在零点不满足任何正指数的 Hölder 估计。
\end{solution}
